\documentclass[12pt,a4paper]{article}
\usepackage[margin=25mm, headheight=15pt, headsep=10pt]{geometry}
\usepackage{fontspec}
\usepackage[table]{xcolor} % Note: table option is required for rowcolors
\usepackage{titlesec}
\usepackage{tocloft}
\usepackage{fancyhdr}
\usepackage{booktabs}
\usepackage{tcolorbox}
\tcbuselibrary{skins, breakable}
\usepackage[colorlinks=true, linkcolor=emerald, urlcolor=emerald, bookmarks=true]{hyperref}

\definecolor{emerald}{RGB}{0,102,51}
\definecolor{darkred}{RGB}{139,0,0}
\definecolor{lightgray}{RGB}{245,245,245}

\usepackage[nil]{babel}
\babelprovide[import, main]{bengali}
\babelprovide[import]{arabic}
\babelfont{rm}[Renderer=HarfBuzz,Scale=1.0]{Noto Serif Bengali}
\babelfont{sf}[Renderer=HarfBuzz]{Noto Sans Bengali}
\babelfont{tt}[Renderer=HarfBuzz]{Noto Sans Bengali}
\babelfont[arabic]{rm}[Renderer=HarfBuzz,Scale=1.1]{Noto Naskh Arabic}

% Section styling
\titleformat{\section}{\Large\bfseries\color{emerald}}{\thesection.}{0.5em}{}
\titleformat{\subsection}{\large\bfseries\color{emerald!85!black}}{\thesubsection.}{0.5em}{}
\titleformat{\subsubsection}{\normalsize\bfseries\color{emerald!70!black}}{\thesubsubsection.}{0.5em}{}
\titlespacing*{\section}{0pt}{18pt}{8pt}

% Fancy Header/Footer setup
\pagestyle{fancy}
\fancyhf{} % clear all header and footer fields
\fancyhead[L]{\color{emerald}\small\textbf{\nouppercase{\leftmark}}}
\fancyfoot[L]{\scriptsize\color{black!50}{\lat{AI}}-সংকলিত, আলেম-যাচাইকৃত নয়}
\fancyfoot[C]{\color{emerald!70!black}\thepage}
\renewcommand{\headrulewidth}{0.4pt}
\renewcommand{\headrule}{\hbox to\headwidth{\color{emerald}\leaders\hrule height \headrulewidth\hfill}}
\renewcommand{\footrulewidth}{0pt}

% Custom boxes
% 1. Ayat/Hadith Box
\newtcolorbox{ayatbox}[1][]{
    enhanced, breakable, 
    colback=emerald!5, 
    colframe=emerald, 
    boxrule=0.5pt, 
    arc=3pt, 
    left=10pt, right=10pt, top=10pt, bottom=10pt,
    #1
}

% 2. Quote/Translation Box
\newtcolorbox{quotebox}[1][]{
    enhanced, breakable,
    frame hidden,
    colback=lightgray,
    borderline west={3pt}{0pt}{emerald},
    left=15pt, right=10pt, top=10pt, bottom=10pt,
    sharp corners,
    #1
}

% 3. AI-generated notice box
\newtcolorbox{warnbox}[1][]{
    enhanced, breakable,
    colback=red!4,
    colframe=darkred,
    boxrule=0.8pt,
    arc=3pt,
    left=10pt, right=10pt, top=10pt, bottom=10pt,
    #1
}

% Commands
\newcommand{\arab}[1]{\foreignlanguage{arabic}{#1}}
\newcommand{\kib}[1]{\textcolor{emerald}{\textbf{#1}}}

% Latin (English) fallback: Noto Serif Bengali has NO Latin glyphs
\newfontfamily\latfont[Renderer=HarfBuzz]{Noto Serif}
\newcommand{\lat}[1]{{\latfont #1}}

\setlength{\parskip}{5pt}
\setlength{\parindent}{0pt}

% Fix for missing bullet in Noto Serif Bengali
\renewcommand{\labelitemi}{$\bullet$}



\begin{document}
\begin{center}{\Large\bfseries\color{emerald} 03-আবু-বকর-ও-উমার-বিবাদে-বিনয়}\end{center}
\vspace{1em}
\section*{ঘটনা ৩: আবু বকর ও উমার (রাঃ) — বিবাদে বিনয়ের শিক্ষা}

\begin{warnbox}
 \textbf{গুরুত্বপূর্ণ নোটিশ (\lat{Important} \lat{Notice}):} এই কন্টেন্টটি \lat{AI} (কৃত্রিম বুদ্ধিমত্তা) দিয়ে প্রস্তুত ও সংকলিত। \lat{AI} কঠোর সূত্র-মান নীতি মেনে শুধু নির্ভরযোগ্য উৎস থেকে তথ্য সংগ্রহ করেছে — কেবল সহীহ/হাসান হাদিস ও সহীহ সনদের আসার রাখা হয়েছে, দুর্বল সনদ বাদ দেওয়া হয়েছে। তবে এটি কোনো আলেম বা মুহাদ্দিস কর্তৃক যাচাইকৃত নয়।
\end{warnbox}


\medskip\hrule\medskip

\subsection*{১. সূত্র কার্ড}

\begin{table}[h]\centering\small
\begin{tabular}{ll}
বিষয় & বিবরণ \\
\hline
\textbf{বর্ণনাকারী} & আবু দারদা (রাঃ) থেকে \\
\textbf{গ্রন্থ} & সহীহ বুখারি, হাদিস ৪৬৪০ \\
\textbf{অধ্যায়} & কিতাবুত তাফসীর (সূরা আ'রাফের প্রেক্ষাপট) \\
\textbf{সনদের মান} & \textbf{সহীহ} \\
\textbf{মূল বিষয়} & ভুল স্বীকার, ক্ষমা চাওয়া ও নবীজির সাথীর প্রতি সম্মান — অহংকারের বিপরীত \\
\hline
\end{tabular}
\end{table}

\medskip\hrule\medskip

\subsection*{২. পটভূমি (কে, কখন, কোথায়, কেন)}

নবীজি (সাল্লাল্লাহু আলাইহি ওয়া সাল্লাম)-এর জীবদ্দশায়, যখন তিনি সাহাবাদের সাথে ছিলেন। \textbf{আবু বকর সিদ্দিক (রাঃ)} — ইসলামের সবচেয়ে প্রিয় ব্যক্তি — এবং \textbf{উমার ইবনুল খাত্তাব (রাঃ)} — দুই মহান সাহাবির মধ্যে একবার \lat{heated} (উত্তপ্ত) \lat{dispute} (বিবাদ) হয়।

দুইজনই দৃঢ় মেজাজের মানুষ ছিলেন। উমার (রাঃ) \lat{anger} (রাগ)-এ চলে গেলেন। আবু বকর (রাঃ) তাঁকে অনুসরণ করে গিয়ে \lat{forgiveness} (ক্ষমা) চাইলেন। কিন্তু রাগের কারণে উমার (রাঃ) তখন ক্ষমা করতে পারেননি — এমনকি দরজা বন্ধ করে দিলেন — একটি \lat{harsh} (কঠোর) প্রতিক্রিয়া।

এই সাধারণ ঘটনা থেকে নবীজি (সা.) একটি \lat{timeless} (চিরন্তন) \lat{lesson} (শিক্ষা) দিলেন — যা এই ঘটনাকে \lat{arrogance} (অহংকার) দমনের পাঠে রূপান্তরিত করেছে।

\medskip\hrule\medskip

\subsection*{৩. পূর্ণ ঘটনার বর্ণনা}

আবু দারদা (রাঃ) বলেন:

\begin{quote}
\textbf{\lat{“}আবু বকর ও উমারের (রাঃ) মধ্যে একবার বিতর্ক হলো। আবু বকর উমারকে রাগিয়ে দিলেন, আর উমার রাগান্বিত হয়ে চলে গেলেন।} \\
\textbf{আবু বকর তাঁর পিছনে গিয়ে বললেন: "আমার জন্য ইস্তিগফার করো" (আমার পক্ষ থেকে আল্লাহর কাছে ক্ষমা চাও/আমাকে ক্ষমা করো)। কিন্তু উমার করলেন না — এমনকি আবু বকরের মুখের সামনে দরজা বন্ধ করে দিলেন।} \\
\textbf{তখন আবু বকর নবীজির (সাল্লাল্লাহু আলাইহি ওয়া সাল্লাম) কাছে এলেন — আমরা তখন তাঁর কাছে বসা ছিলাম। নবীজি বললেন: "তোমাদের এই সাথী নিশ্চয়ই কারো সাথে ঝগড়া করেছে।"} \\
\textbf{এদিকে উমার (রাঃ) যা করেছিলেন তাতে অনুতপ্ত হলেন। তিনি এসে সালাম দিয়ে নবীজির কাছে বসলেন এবং ঘটনাটি খুলে বললেন। তখন নবীজি (সা.) রাগান্বিত হলেন — আর আবু বকর বলতে লাগলেন: "হে আল্লাহর রাসূল! আল্লাহর কসম! এর (উমারের) চেয়ে বেশি অন্যায় আমি করেছি।"} \\
\textbf{নবীজি (সা.) বললেন: "তোমরা কি আমার সাথীকে (আবু বকরকে) আমার জন্য ছেড়ে দেবে না? তোমরা কি আমার সাথীকে আমার জন্য ছেড়ে দেবে না? আমি যখন বললাম — 'হে মানুষ! আমি তোমাদের সবার কাছে আল্লাহর রাসূল', তখন তোমরা বলেছিলে — 'তুমি মিথ্যা বলছ', কিন্তু আবু বকর বলেছিলেন — 'আপনি সত্য বলেছেন!'"\lat{”}}
\end{quote}

আবু বকর (রাঃ) নিজে নবীজির পায়ের কাছে ঝুঁকে পড়ে বললেন: \textbf{\lat{“}হে আল্লাহর রাসূল! আল্লাহর কসম! এর (উমারের) চেয়ে বেশি অন্যায় আমি করেছি।\lat{”}} — নবীজি দুবার বললেন: \textbf{\lat{“}তোমরা কি আমার সাথীকে কষ্ট দেওয়া বন্ধ করবে না?\lat{”}} এরপর কেউ আর আবু বকরকে কষ্ট দেয়নি।

\medskip\hrule\medskip

\subsection*{৪. ধাপে ধাপে বিশ্লেষণ}

\textbf{ধাপ ১ — বিতর্ক ও রাগ:} দুই সাহাবির মধ্যে মতপার্থক্য; উমার \lat{anger} (রাগ)-এ চলে যান — একটি \lat{impulsive} (আবেগপ্রবণ) \lat{reaction} (প্রতিক্রিয়া)।

\textbf{ধাপ ২ — আবু বকরের \lat{humility} (বিনয়):} আবু বকর (রাঃ) যিনি বয়স ও মর্যাদায় উমারের চেয়ে বড়, তিনি নিজে গিয়ে ক্ষমা চান। এটি \lat{arrogance} (অহংকার)-এর সম্পূর্ণ বিপরীত — নিজেকে ছোট করে \lat{forgiveness} (ক্ষমা) চাওয়া।

\textbf{ধাপ ৩ — উমারের দৃঢ়তা:} \lat{anger} (রাগ)-এর মাথায় উমার (রাঃ) ক্ষমা করতে পারেননি। এটাও মানবিক \lat{weakness} (দুর্বলতা) — কিন্তু এর সমাধানও ঘটনায় আছে।

\textbf{ধাপ ৪ — নবীজির সিদ্ধান্ত:} নবীজি (সা.) উমারের অনুতাপের পর দুইজনকেই শান্ত করতে চাইলেন — একটি \lat{reconciliation} (সমঝোতা) চেষ্টা। আবু বকর বারবার নিজের দোষ স্বীকার করলেন — \lat{“}আমি বেশি অন্যায় করেছি\lat{”} — এটাই সাহাবাদের \lat{character} (চরিত্র)-এর উচ্চতা।

\textbf{ধাপ ৫ — নবীজির চূড়ান্ত \lat{lesson} (শিক্ষা):} নবীজি আবু বকরের \lat{dignity} (মর্যাদা) তুলে ধরলেন — তিনি ইসলামের প্রথম সাক্ষী, মক্কার সবচেয়ে কঠিন মুহূর্তে নবীজির সাথে থাকা সাথী। নবীজি বললেন: \lat{“}আমার সাথীকে ছেড়ে দাও।\lat{”} — ফলে কেউ আর আবু বকরকে কষ্ট দেওয়ার সাহস পায়নি।

\medskip\hrule\medskip

\subsection*{৫. শব্দের ব্যাখ্যা}

\begin{table}[h]\centering\small
\begin{tabular}{ll}
শব্দ & অর্থ \\
\hline
\textbf{মুহাওয়ারাহ (\arab{محاورة})} & কথোপকথন/বিতর্ক — এখানে মতপার্থক্যের কারণে উত্তপ্ত বাক্য বিনিময় \\
\textbf{আওজাউ (\arab{أوذا})} & দুঃখ/কষ্ট দেওয়া \\
\textbf{সাহিবি (\arab{صاحبي})} & আমার সাথী — নবীজির আবু বকরের প্রতি স্নেহভরা সম্বোধন \\
\textbf{ইস্তিগফার} & ক্ষমা প্রার্থনা — আবু বকর উমারকে নিজের পক্ষে ক্ষমা চাওয়ার কথা বলেন \\
\hline
\end{tabular}
\end{table}

\textbf{মূল বাক্য — \lat{“}আমি বেশি অন্যায় করেছি\lat{”}:}
আবু বকর (রাঃ) নবীজির সামনে ও উমারের উপস্থিতিতে বারবার নিজের দোষ স্বীকার করেছেন। অহংকারী ব্যক্তি কখনো নিজের দোষ স্বীকার করে না; বিনয়ী ব্যক্তিই এই কাজ করে।

\medskip\hrule\medskip

\subsection*{৬. সংশ্লিষ্ট হাদিস ও আয়াত}

\textbf{হাদিস (বুখারি ৩৬৬৮):} নবীজি (সা.) বলেছেন: \textbf{\lat{“}আমি যদি মানুষকে বন্ধু (খলিল) বানাতাম, তবে আবু বকরকে বানাতাম।\lat{”}} — আবু বকরের মর্যাদা।

\textbf{হাদিস (মুসলিম ৯১):} নবীজি বলেছেন: \textbf{\lat{“}কিবর হলো সত্যকে তুচ্ছ করা ও মানুষকে হেয় করা।\lat{”}} — আবু বকর সত্যকে তুচ্ছ করেননি; বরং নিজের দোষ স্বীকার করেছেন, এটাই বিনয়।

\textbf{আয়াত (সূরা আলে ইমরান ৩:১৩৪):}
\begin{quote}
\textbf{\lat{“}যারা সুখে-দুঃখে দান করে, ক্রোধ সংবরণ করে এবং মানুষকে ক্ষমা করে — আল্লাহ সৎকর্মশীলদের ভালোবাসেন।\lat{”}}
\end{quote}

\textbf{আয়াত (সূরা আলে ইমরান ৩:১৫৯):}
\begin{quote}
\textbf{\lat{“}আল্লাহর রহমতেই আপনি তাদের প্রতি কোমল ছিলেন। আপনি যদি রূঢ় ও কঠোর হৃদয় হতেন, তবে তারা আপনার চারপাশ থেকে ছত্রভঙ্গ হয়ে যেত।\lat{”}}
\end{quote}

\medskip\hrule\medskip

\subsection*{৭. গভীর শিক্ষা}

\textbf{১. ভুল স্বীকার করা \lat{arrogance} (অহংকার) ধ্বংসের প্রথম ধাপ।}
আবু বকর (রাঃ) — ইসলামের সর্বশ্রেষ্ঠ মানুষ — তিনিও মানুষের মতো ভুল করেছেন, কিন্তু ভুলের পর তাঁর \lat{reaction} (প্রতিক্রিয়া) ছিল দ্রুত ক্ষমা চাওয়া। \lat{Arrogant} (অহংকারী) ভুল লুকায়, \lat{humble} (বিনয়ী) মেনে নেয়।

\textbf{২. \lat{anger} (রাগ)-এ শক্ত কথা — পরে অনুতাপ।}
উমার (রাঃ) রাগে ক্ষমা করতে পারেননি, কিন্তু শীঘ্রই অনুতপ্ত হয়ে ফিরে এসেছেন — একটি \lat{positive} (ইতিবাচক) \lat{step} (পদক্ষেপ)। রাগ স্বাভাবিক; কিন্তু স্থায়ী \lat{arrogance} (অহংকার)-এ থেকে যাওয়াই ধ্বংসের কারণ। উমার অনুতপ্ত হয়ে ফিরে আসায় নবীজিও তাকে শান্ত করেছেন।

\textbf{৩. নিজের চেয়ে অন্যজনের ভুল স্বীকার করা — সাহাবাদের \lat{greatness} (মহত্ত্ব)।}
উমার আসার পর আবু বকর বলেছেন: \lat{“}আমিই বেশি অন্যায় করেছি।\lat{”} — নিজেকে ছোট করে বলা, অপরকে বাঁচানো — এটাই \lat{humility} (বিনয়)-এর চরম শিখর।

\textbf{৪. নবীজির সাথীদের প্রতি \lat{respect} (সম্মান) রক্ষা।}
নবীজি (সা.) আবু বকরের মর্যাদা রক্ষায় স্পষ্ট ভাষায় বলেছেন: \lat{“}আমার সাথীকে ছেড়ে দাও।\lat{”} — \lat{lesson} (শিক্ষা): যারা ইসলামের \lat{service} (সেবা)-এ মহান, তাদের সম্মান করা ঈমানের অংশ; আর তাদের সম্মানে অহংকার করা কিবর।

\textbf{৫. অহংকারী ও বিনয়ীর \lat{result} (ফলাফল)।}
আবু বকর বিনয়ের কারণে জান্নাতের সুসংবাদ পেয়েছেন (বুখারি ৩৬৬৮); অহংকারীর জন্য আছে জাহান্নামের \lat{punishment} (শাস্তি)। \lat{humility} মানুষকে বড় করে, \lat{arrogance} ছোট করে।

\medskip\hrule\medskip

\subsection*{৮. জীবনে প্রয়োগের পদ্ধতি}

\begin{enumerate}
\item \textbf{ভুল হলে সাথে সাথে স্বীকার করুন:} কারো সাথে বিতর্কে নিজের ভুল বুঝলে আবু বকরের মতো \lat{“}আমিই ভুল করেছি\lat{”} বলা শিখুন — এটাই \lat{integrity} (সততা)।
\item \textbf{রাগ \lat{control} (নিয়ন্ত্রণ):} রাগ হলে চুপ থাকুন, পরে শান্ত মনে সমাধান করুন — উমারের অনুতাপের পথ অনুসরণ করুন।
\item \textbf{ক্ষমা চাইতে দেরি করবেন না:} আবু বকর রাগের পরপরই ক্ষমা চেয়েছেন; দেরি করলে \lat{arrogance} (অহংকার) বেড়ে যায়।
\item \textbf{নবীজির সাথীদের ভালোবাসা:} সাহাবাদের প্রতি ভালোবাসা ঈমানের অংশ (তিরমিজি) — তাদের জীবন থেকে \lat{lesson} (শিক্ষা) নিন।
\item \textbf{বিবাদে নিজের দোষ খোঁজুন:} অপরকে \lat{blame} (দোষারোপ) না করে নিজের ভাগের দোষ খুঁজে স্বীকার করুন।
\end{enumerate}
\medskip\hrule\medskip

\subsection*{৯. রেফারেন্স}

\begin{itemize}
\item সহীহ বুখারি, হাদিস ৪৬৪০ (কিতাবুত তাফসীর)
\item সহীহ বুখারি, হাদিস ৩৬৬৮ (আবু বকরের মর্যাদা)
\item সহীহ মুসলিম, হাদিস ৯১ (কিবরের সংজ্ঞা)
\item সূরা আলে ইমরান, আয়াত ১৩৪ ও ১৫৯
\end{itemize}

\end{document}
